% Zakljucak

У овом мастер раду разматрана је проблематика ефикасне обраде и управљања динамичким графовима, који представљају кључни математички и рачунарски модел за репрезентацију савремених комплексних система склонх променама током времена. Традиционални приступи, који се ослањају на анализаторе статичких графова или периодично поновно израчунавање свих својстава система, показују озбиљна ограничења у погледу временске и просторне сложености када се примене на окружења са високом фреквенцијом ажурирања.

Кроз детаљну анализу постојећих имплементационих модела из литературе, у раду је истакнута потреба за структуром података која балансира између брзине структурних модификација и ефикасности аналитичких пролаза кроз суседства чворова. Тежиште рада стављено је на структуру хеш-индексираних блокова суседства. DHB структура кроз комбинацију секвенцијалног смештања суседа у контигуалне блокове и хеш-индекса са отвореним адресирањем успешно спаја добре особине низова (константан индексни приступ, просторна локалност, ефикасна итерација) и хеш-табела (брза провера постојања гране и промена тежине у $O(1)$ времену).

Главни истраживачки и практични допринос овог рада огледа се у реализацији C++ библиотеке засноване на DHB структури и њеној успешној адаптацији за решавање фундаменталних графовских проблема у динамичким условима:
\begin{itemize}
    \item \textit{Инкрементална динамичка повезаност:} Интеграцијом DHB структуре са помоћном структуром дисјунктних подскупова омогућено је праћење компоненти повезаности и одговарање на упите о повезаности чворова у амортизованом готово константном времену, уз задржавање ефикасних инкременталних измена.
    \item \textit{Динамичко минимално разапињуће стабло:} Искоришћена је флексибилност DHB структуре у одржавању произвољног и сортираног поретка суседа. Применом мин-хипа над унапред сортираним блоковима суседства елиминисана је потреба за скупим глобалним сортирањем свих грана графа, чиме је постигнута значајно боља асимптотска ефикасност одржавања минималне разапињуће шуме након сваке промене тежина или структуре.
    \item \textit{Динамичко максимално тежинско упаривање:} Приказано је како одржавање опадајућег поретка тежина у блоковима суседства директно погодује похлепним алгоритмима апроксимације. DHB омогућава дохватање гране са највећом тежином у $O(1)$ времену преко реверзних итератора, што омогућава брзу евалуацију понуда у Суитор алгоритму и ефикасно одржавање $1/2$-максималног тежинског упаривања.
\end{itemize}

Иако разматрани алгоритми нису имплементирани потпуно динамички, већ се одређена количина информација израчунава изнова, примарни фокус био је на демонстрацији особина DHB структуре у погледу ефикасности и могућности прилагођавања динамичким променама. Важно је нагласити да се DHB структура не може користити као потпуно самостално решење за све динамичке алгоритме, нити сама по себи покрива све функције неопходне за различите примене. Уместо тога, она служи као изузетно ефикасно складиште основних градивних елемената графа. Тиме је показано да DHB структура представља солидну основу за развој потпуно динамичких алгоритама који би омогућили одржавање својстава графа у присутству и инкременталних и декременталних операција.

Поред теоријске и имплементационе анализе, у раду је кроз репрезентативне примере показано да динамички графови нису само апстрактни концепт, већ да постоји реална и све већа потреба за њиховом применом, како у научним тако и у индустријским окружењима.

Као главни правци будућег истраживања и надградње овог рада истичу се проширење DHB структуре ради подршке за потпуно динамичке алгоритме са ефикасним ажурирањима, као и имплементација напредних паралелних и дистрибуираних варијанти ове структуре. Такође, значајан фокус будућег рада може бити на даљој оптимизацији управљања меморијом приликом фреквентних операција реалокације и поновног хеширања, применици овог модела у динамичком машинском учењу над графовима (енгл. \textit{Dynamic Graph Neural Networks}), као и у развоју графовских база података са високом фреквенцијом измена где динамичке структуре играју кључну улогу у ефикасном управљању и анализи података.

% Zakljucak